\chapter{Clase 3}
\section{Estructura y Fuentes de los Datos}

\subsection{Datos}

Los datos son las observaciones recolectadas (como mediciones, géneros, respuestas de encuesta). Se recogen, analizan y resumen para su presentación e interpretación. A todos los datos reunidos para un determinado estudio se les llama \textbf{conjunto de datos} para el estudio.

\begin{longtblr}{colspec={|*{8}{X[c]}|}, rowhead=1, cells={font=\small}}
\hline
\textbf{Sexo} & \textbf{Nivel de estudio} & \textbf{Años de estudio} & \textbf{Medio pref.} & \textbf{Ingresos (S.M)} & \textbf{Edad} & \textbf{Estatura} & \textbf{Estrato} \\
\hline
F & Pregrado & 16 & Prensa & 1.0 & 23 & 1.64 & 3 \\
\hline
M & Bachiller & 13 & TV & 1.5 & 23 & 1.65 & 4 \\
\hline
M & Bachiller & 12 & Radio & 1.0 & 23 & 1.60 & 3 \\
\hline
M & Pregrado & 18 & Prensa & 2.5 & 51 & 1.57 & 3 \\
\hline
M & Media-Básica & 10 & TV & 1.0 & 19 & 1.77 & 2 \\
\hline
\end{longtblr}

\subsection{Fuentes de Datos}

Las fuentes de datos pueden clasificarse de la siguiente manera:

\begin{itemize}
    \item \textbf{Registros rutinarios:} Hospitales, registros contables, entre otros.
    \item \textbf{Encuestas:} Se recurre a ellas cuando no hay registros almacenados.
    \item \textbf{Experimentación:} Resultados de un diseño experimental.
    \item \textbf{Fuentes externas:} Publicaciones de instituciones.
\end{itemize}

En ocasiones, los datos no pueden obtenerse de las fuentes existentes y deben generarse mediante la realización de un estudio estadístico. Estos datos se clasifican como \textbf{experimentales} u \textbf{observacionales}.

\subsection{Variable}

Una \textbf{variable} es una característica de los elementos que es de interés para el estudio. Un conjunto de datos puede contener múltiples variables. Por ejemplo, en la tabla anterior se identifican variables como sexo, nivel de estudio, años de estudio, medio de comunicación preferido, ingresos, edad, estatura y estrato socioeconómico.

\subsection{Elementos}

Los \textbf{elementos} son las entidades de las que se obtienen los datos. En un conjunto de datos, los nombres o identificadores de los elementos aparecen como referencia de cada observación registrada.

\section{Datos de Sección Transversal y de Series de Tiempo}

Para los propósitos del análisis estadístico, la distinción entre datos transversales y datos de series de tiempo es importante.

\subsection{Datos de Sección Transversal}

Los datos de sección transversal son los obtenidos en el mismo o aproximadamente el mismo momento (punto en el tiempo). Por ejemplo, una tabla que describe múltiples variables de varios individuos en un mismo instante constituye un conjunto de datos transversales.

\subsection{Datos de Series de Tiempo}

Los datos de series de tiempo son datos obtenidos a lo largo de varios periodos. Los economistas, por ejemplo, utilizan este tipo de datos para identificar tendencias.

\section{Errores en la Adquisición de Datos}

Usar datos erróneos es peor que no usar ningún dato. Un \textbf{error} ocurre cuando el dato obtenido no es igual al verdadero valor.

\subsection{Tipos de Errores}

\begin{itemize}
    \item Error de escritura: registrar 24 años en lugar de 42 años.
    \item Malinterpretación de la pregunta por parte del entrevistado, lo que genera una respuesta incorrecta.
\end{itemize}

Se debe tener sumo cuidado tanto al recolectar los datos como al registrarlos para garantizar que no se cometan errores.

\subsection{Verificación de Consistencia}

Es necesario verificar la consistencia interna de los datos. Por ejemplo, un entrevistado con 24 años no puede tener 20 años de experiencia laboral. Se deben revisar datos que tengan valores inusualmente grandes o pequeños, denominados \textbf{observaciones atípicas}, que son candidatos a posibles errores.

\section{Razones Prácticas para Utilizar una Muestra}

Aunque un censo proporciona información de todos los miembros de una población, no siempre es viable realizarlo. El uso de una muestra puede reducir los costos y el tiempo requerido para el estudio, además de permitir un mayor control durante la obtención de los datos. En algunas situaciones, el muestreo es la única alternativa práctica debido a que:

\begin{itemize}
    \item La población es tan grande que excede las posibilidades del investigador, tanto económicas como por limitaciones temporales.
    \item La población es suficientemente homogénea para estudiar una fracción de ella sin censar todos sus elementos.
    \item El proceso de medida o investigación es destructivo, como ocurre al consumir un cierto artículo para juzgar su calidad.
\end{itemize}

La principal limitación es la variabilidad muestral: dos muestras obtenidas de una misma población pueden producir resultados diferentes. Por esta razón, el alcance de las conclusiones debe considerar el tamaño de la muestra y la precisión requerida por el estudio.

\section{Pensamiento Crítico ante la Estadística}

\subsection{Usos Inadecuados de la Estadística}

Por lo general, existen dos situaciones en que la estadística se utiliza como fuente de engaño:

\begin{enumerate}
    \item El intento malintencionado por parte de individuos deshonestos.
    \item Errores no intencionales por parte de personas que no dominan la disciplina.
\end{enumerate}

\subsection{Gráficas y su Uso Inadecuado}

Las gráficas pueden distorsionar la percepción de los datos si no se construyen correctamente. Se debe verificar que las escalas sean adecuadas y que no se exageren diferencias.

\subsection{Resultados Reportados}

Cuando se obtienen datos de personas, es mejor tomar las medidas directamente en lugar de pedir a los sujetos que reporten resultados. Las personas tienden a redondear, generalmente hacia abajo, y en temas como ingresos tienden a mentir.

\subsection{Otros Aspectos a Considerar}

\begin{itemize}
    \item Uso adecuado de porcentajes.
    \item Preguntas que inducen respuestas sesgadas.
    \item Orden de las preguntas en una encuesta.
    \item Falta de respuesta por parte de los encuestados.
    \item Datos faltantes.
    \item Estudios realizados para el propio beneficio del investigador.
    \item Números precisos que pueden ser falsos.
    \item Distorsiones deliberadas.
\end{itemize}

\section{Recolección de Datos Muestrales}

\subsection{Fundamentos de la Recolección de Datos}

Los métodos estadísticos se rigen por los datos recabados. Por lo regular, se obtienen datos de dos fuentes distintas:

\begin{longtblr}{colspec={|Q[l,wd=3.5cm]|Q[l,wd=9cm]|}}
\hline
\textbf{Estudio observacional} & Se observan y miden características específicas, pero no se intenta modificar a los sujetos que se están estudiando. \\
\hline
\textbf{Experimento} & Se aplican algunos tratamientos y luego se procede a observar sus efectos sobre los sujetos. Los sujetos se denominan \textit{unidades experimentales}. \\
\hline
\end{longtblr}

\subsection{Estudios Retrospectivos y Prospectivos}

\begin{longtblr}{colspec={|Q[l,wd=3.5cm]|Q[l,wd=9cm]|}}
\hline
\textbf{Estudio retrospectivo (de control de caso)} & Los datos se toman del pasado (mediante examen de registros, entrevistas y otros recursos). \\
\hline
\textbf{Estudio prospectivo (longitudinal o de cohorte)} & Los datos se reunirán en el futuro y se toman de grupos (cohortes) que comparten factores comunes. \\
\hline
\end{longtblr}

\section{Estructura de un Conjunto de Datos}

Un conjunto de datos puede organizarse de diversas formas. Una estructura común es la tabla bidireccional, donde las filas representan elementos (o categorías) y las columnas representan variables. Por ejemplo:

\begin{longtblr}{colspec={|*{4}{X[c]}|}, rowhead=1, cells={font=\small}}
\hline
 & \textbf{Prensa} & \textbf{Radio} & \textbf{TV} \\
\hline
Varón & 30 & 40 & 90 \\
\hline
Mujer & 25 & 35 & 75 \\
\hline
\end{longtblr}

De manera más general, si se tienen $i$ niveles de estudio, $j$ medios de comunicación preferido y $k$ categorías de sexo, los datos se pueden organizar en una estructura tridimensional donde cada celda $(i, j, k)$ contiene la frecuencia correspondiente.


\section{Pensamiento Estadístico}

\subsection{Concepto Clave}

Al finalizar el estudio introductorio de la estadística, el estudiante cuenta con numerosas herramientas analíticas. Sin embargo, efectuar cálculos sin considerar aspectos generales importantes equivale a estar ``equipado peligrosamente''. El pensamiento estadístico exige que, antes de cualquier análisis formal, se tomen en cuenta factores fundamentales que determinan la validez y utilidad de los resultados.

Para aprender a pensar en términos estadísticos, suelen ser más importantes el sentido común y las consideraciones prácticas que la aplicación irreflexiva de fórmulas y cálculos. Para analizar de forma adecuada unos datos, es necesario contar con información adicional que proporcione el contexto adecuado.

Las preguntas fundamentales que guían el pensamiento estadístico son:

\begin{itemize}
    \item ¿Cuál es el contexto de los datos?
    \item ¿De qué fuente se obtuvieron?
    \item ¿Cómo se recabaron?
    \item ¿Qué se puede concluir a partir de la información?
    \item Con base en conclusiones estadísticas, ¿qué implicaciones prácticas resultan del análisis?
\end{itemize}

\subsection{Contexto y Fuente de los Datos}

Siempre se debe tomar en cuenta el contexto de los datos, ya que este determina el análisis estadístico que debe emplearse. El contexto incluye la naturaleza de las variables, las unidades de medición y el propósito del estudio.

Respecto a la fuente, es necesario considerar si esta es objetiva o si existe alguna razón para pensar que está sesgada. No todos los estudios cuentan con fuentes sin sesgo. Se debe permanecer atento y escéptico ante estudios que provienen de fuentes que podrían tener intereses particulares en obtener ciertos resultados.

\subsection{Método de Muestreo}

Al reunir datos muestrales para un estudio, el método de muestreo elegido puede afectar de manera importante la validez de las conclusiones. Un aspecto crítico es que las muestras de respuesta voluntaria (o autoseleccionadas) a menudo están sesgadas, ya que es más probable que los individuos con un interés especial en el tema decidan participar. En una muestra de respuesta voluntaria, los propios sujetos deciden participar, y aunque es posible utilizar métodos estadísticos válidos para analizarlas, los resultados no son necesariamente válidos.

\subsection{Conclusiones e Implicaciones Prácticas}

Al obtener conclusiones a partir de un análisis estadístico, es necesario formular afirmaciones que sean claras para personas sin conocimientos de estadística ni de su terminología. Se debe evitar cuidadosamente realizar afirmaciones que no estén justificadas por el análisis estadístico.

Además de plantear conclusiones claras, también se deben identificar las implicaciones prácticas de los resultados. Una conclusión puede ser estadísticamente correcta pero carecer de relevancia práctica, o viceversa.

\subsection{Significancia Estadística y Significancia Práctica}

La significancia estadística de un estudio difiere de su significancia práctica. Es posible que, con base en los datos muestrales disponibles, se utilicen métodos estadísticos para llegar a la conclusión de que algún tratamiento o hallazgo es eficaz, aunque el sentido común sugiera que la diferencia no es suficiente para justificar su uso práctico.

En términos formales, la significancia estadística se alcanza cuando la probabilidad de obtener el resultado observado por puro azar es muy pequeña. Un criterio común es que se logra significancia estadística si dicha probabilidad es del $5\%$ o menos:

$$P(\text{resultado por azar}) \leq 0.05$$

Sin embargo, un resultado estadísticamente significativo no necesariamente tiene significancia práctica. Por ejemplo, si un tratamiento aumenta la probabilidad de un evento del $50\%$ al $52\%$, el resultado puede ser estadísticamente significativo con una muestra suficientemente grande, pero la diferencia práctica es mínima.

\subsection{Definición de Pensamiento Estadístico}

El pensamiento estadístico incluye el pensamiento crítico y la capacidad de interpretar los resultados. También puede implicar determinar si los resultados son estadísticamente significativos. El pensamiento estadístico va mucho más allá de la simple capacidad de ejecutar cálculos complejos; implica la capacidad de observar el panorama general, tomar en cuenta factores relevantes como el contexto, la fuente de los datos y el método de muestreo, obtener conclusiones e identificar implicaciones prácticas.

\section{Tipos de Datos}

\subsection{Parámetro y Estadístico}

Es importante distinguir entre parámetro y estadístico. Los primeros son cantidades constantes que describen una población; los segundos dependen de la muestra extraída y varían de una muestra a otra.

\begin{itemize}
    \item \textbf{Parámetro:} es cualquier cantidad constante que describe una característica de la población (por ejemplo, la media poblacional $\mu$ o la proporción poblacional $p$).
    \item \textbf{Estadístico:} es cualquier función de las variables aleatorias que constituyen una muestra extraída de la población (por ejemplo, la media muestral $\bar{x}$ o la proporción muestral $\hat{p}$).
\end{itemize}

La notación habitual asocia parámetros con letras griegas y estadísticos con letras latinas:

$$\text{Parámetros poblacionales: } \mu,\; \sigma,\; p \qquad \text{Estadísticos muestrales: } \bar{x},\; s,\; \hat{p}$$

\subsection{Datos Cuantitativos}

Los datos cuantitativos (o numéricos) consisten en números que representan conteos o mediciones. Se subdividen en dos categorías:

\begin{itemize}
    \item \textbf{Datos discretos:} resultan cuando el número de posibles valores es finito o contable. Suelen tomar valores enteros y existe separación o interrupción entre valores consecutivos. Ejemplos: número de hijos, cantidad de automóviles, número de defectos en un lote.
    \item \textbf{Datos continuos:} resultan de un infinito de posibles valores que pueden asociarse a puntos de una escala continua, cubriendo un rango sin huecos ni interrupciones. Pueden tomar cualquier valor dentro de un intervalo. Ejemplos: estatura, peso, temperatura, tiempo de reacción.
\end{itemize}

\subsection{Datos Cualitativos}

Los datos cualitativos (también llamados categóricos o de atributo) se dividen en diferentes categorías que se distinguen por alguna característica no numérica. No se pueden medir numéricamente de forma directa. Ejemplos: nivel de estudios, género, religión, raza, estrato socioeconómico, nacionalidad, color de la piel.

\subsection{Niveles de Medición}

Otra forma de clasificar los datos consiste en usar cuatro niveles de medición. Cuando se aplica la estadística a problemas reales, el nivel de medición de los datos es un factor importante para determinar el procedimiento estadístico adecuado.

\textbf{Nivel nominal.} Los datos consisten exclusivamente en nombres, etiquetas o categorías cuyos valores solo poseen la propiedad de identidad ($a = a$, $a \neq b$). No se pueden organizar en un esquema de orden ni realizar operaciones aritméticas con ellos. Ejemplo: color de ojos, género, ciudad de residencia.

\textbf{Nivel ordinal.} Los datos pueden acomodarse en algún orden o jerarquía, existe transitividad ($a < b$ y $b < c$ implica $a < c$), aunque no es posible determinar diferencias entre los valores o tales diferencias carecen de significado. Ejemplo: clasificaciones de universidades, nivel de satisfacción (bajo, medio, alto).

\textbf{Nivel de intervalo.} En contraste con el nivel ordinal, esta escala tiene diferencias precisas entre unidades de medida. Establece un orden en las posiciones relativas y la diferencia entre dos valores cualesquiera tiene significado. Sin embargo, no posee un punto de partida cero inherente (natural); el cero es relativo y arbitrario, por lo que las razones no tienen sentido. Por ejemplo, una temperatura de $0^\circ\text{C}$ no indica ausencia de temperatura, y no es válido afirmar que $50^\circ\text{F}$ es el doble de temperatura que $25^\circ\text{F}$. Un coeficiente intelectual de $0$ no implica ausencia de inteligencia.

\textbf{Nivel de razón.} Posee todas las características del nivel de intervalo, con la propiedad adicional de un punto de partida cero inherente donde cero indica que nada de la cantidad está presente. Para valores en este nivel, tanto las diferencias como las proporciones tienen significado. Ejemplos: peso ($60\;\text{kg}$ es un $20\%$ más pesado que $50\;\text{kg}$), distancia, tiempo, edad. El cero absoluto o natural representa la nulidad de lo que se estudia.

Una regla práctica para distinguir entre nivel de intervalo y nivel de razón es considerar si la expresión ``dos veces'' tiene sentido: si una cantidad es el doble de la otra y esta afirmación es válida, se aplica el nivel de razón; en caso contrario, se aplica el nivel de intervalo.